
\paragraph{\textbf{Corollary 3.2.}}
\SGARefTarget{sga3:xiv:3.2}
{\itshape Let \(G\) and \(\mathfrak g\) be as in \SGARefLink{sga3:xiv:3.1}{3.1}, with \(G\)
of finite type over \(S\), and let \(\mathfrak k\) and
\(\mathfrak h\) be two Lie subalgebras of \(\mathfrak g\) that are
locally direct factors as modules. Suppose that one of the following
two hypotheses holds:

\smallskip
\noindent
\textup{a)}
For every \(s\in S\), the geometric fiber
\(\mathfrak k_{\bar s}\) is nilpotent and contains a regular element
of \(\mathfrak g_{\bar s}\), while the geometric fiber
\(\mathfrak h_{\bar s}\) contains a Cartan subalgebra of
\(\mathfrak g_{\bar s}\).

\smallskip
\noindent
\textup{b)}
\(\mathfrak k\) is a Cartan subalgebra of \(\mathfrak g\).

\smallskip
\noindent
Then \(\operatorname{Transp}_G(\mathfrak k,\mathfrak h)\) is a
closed subprescheme of \(G\), smooth over \(S\). Moreover, in case
\textup{a)}, its structure morphism to \(S\) is surjective.}

The representability of the transporter by a closed subprescheme of
\(G\) of finite presentation is immediate and is left to the
reader. We first prove smoothness in case \textup{a)}. Suppose first
that \(\mathfrak k\) has a section \(a\) that is quasi-regular in
\(\mathfrak g\). It is then enough to apply \SGARefLink{sga3:xiv:3.1}{3.1} and the following
lemma.

\paragraph{\textbf{Lemma 3.3.}}
\SGARefTarget{sga3:xiv:3.3}
{\itshape Under the hypotheses of \SGARefLink{sga3:xiv:3.2}{3.2}\textup{(a)}, if \(a\) is a
section of \(\mathfrak k\) that is quasi-regular in
\(\mathfrak g\), then
\[
  \operatorname{Transp}_G(\mathfrak k,\mathfrak h)
  =
  \operatorname{Transp}_G(a,\mathfrak h).
\]}

Indeed, by the definitions it is enough to show that, if \(a\) is
also a section of \(\mathfrak h\), then
\(\mathfrak k\subset\mathfrak h\). Since \(\mathfrak k\) is
locally nilpotent by hypothesis, \SGARefLink{sga3:xiv:2.1}{2.1} gives
\[
  \mathfrak k\subset\Nil(a,\mathfrak g).
\]
On the other hand,
\[
  \Nil(a,\mathfrak g)\subset\mathfrak h,
\]
because
\[
  \operatorname{ad}(a)_{\mathfrak g/\mathfrak h}
\]
is injective; this can be checked fiber by fiber by \SGARefLink{sga3:xiii:proposition:4:8:item:b}{XIII~4.8(b)}.
The lemma follows.

In the general case, the usual standard argument first reduces us
to the case where \(S\) is affine noetherian, and then to the case
where \(S\) is local artinian, smoothness being an infinitesimal
property. By flat descent we may further suppose that the residue
field is infinite. The fiber of \(\mathfrak k\) then contains an
element regular in the corresponding fiber of \(\mathfrak g\).
Lifting this element to a section of \(\mathfrak k\) reduces us to
the preceding case.

Thus the transporter is smooth in case \textup{a)}. Its structure
morphism is surjective because this may be checked after reducing to
the case where \(S\) is the spectrum of an algebraically closed
field; \(\mathfrak k\) then contains a regular point of
\(\mathfrak g\), and \SGARefLink{sga3:xiv:3.3}{3.3} and \SGARefLink{sga3:xiv:3.1}{3.1} apply.

For case \textup{b)}, the definition of smoothness \textup{(\SGARefLink{sga3:xi:definition:1:1}{XI~1.1})}
reduces the assertion to the following lifting statement. Let \(S\)
be affine, let \(S_0\) be a closed subscheme defined by a
quasi-coherent nilpotent ideal \(J\), and let
\(g_0\in G(S_0)\) transport \(\mathfrak k_0\) into
\(\mathfrak h_0\). Then \(g_0\) lifts to an element \(g\in G(S)\)
that transports \(\mathfrak k\) into \(\mathfrak h\). The
hypothesis on \(g_0\) places this lifting problem under case
\textup{a)}, already proved. This completes the proof.

Of course, if in \SGARefLink{sga3:xiv:3.2}{3.2}\textup{(b)} the subalgebra \(\mathfrak h\)
satisfies the stronger hypothesis of \SGARefLink{sga3:xiv:3.1}{3.1}, then, and only then, the
structure morphism
\[
  \operatorname{Transp}_G(\mathfrak k,\mathfrak h)
  \longrightarrow S
\]
is surjective. Using Hensel's lemma \SGARefLink{sga3:xi:proposition:1:10}{XI~1.10}, one obtains from \SGARefLink{sga3:xiv:3.1}{3.1}
and \SGARefLink{sga3:xiv:3.2}{3.2} the following corollary.
